\section{Experiment: Standard Library and Syscalls}

\subsection{Objective}
Explore the relationship between C standard library functions and system calls.

\subsection{Setup Requirements}
\begin{itemize}
\item GCC compiler
\item strace tool
\item Linux environment
\end{itemize}

\subsection{Part 1: Tracing System Calls}

Trace a simple program:

\begin{lstlisting}[language=bash]
strace ./simple_program
strace -c ./simple_program
\end{lstlisting}

\subsection{Part 2: Direct Syscalls}

Make system calls directly:

\begin{lstlisting}[language=C]
#include <unistd.h>
#include <sys/syscall.h>

int main(void) {
    // Using syscall wrapper
    syscall(SYS_write, 1, "Hello\n", 6);
    
    // Using inline assembly (x86-64)
    asm volatile (
        "mov $1, %%rax\n"
        "mov $1, %%rdi\n"
        "syscall\n"
        : : "S"("Hello\n"), "d"(6)
        : "rax", "rdi"
    );
    return 0;
}
\end{lstlisting}

\subsection{Part 3: stdio vs write}

Compare buffered I/O with direct writes:

\begin{lstlisting}[language=C]
#include <stdio.h>
#include <unistd.h>

int main(void) {
    // stdio (buffered)
    printf("Buffered output\n");
    
    // Direct syscall (unbuffered)
    write(1, "Direct output\n", 14);
    
    return 0;
}
\end{lstlisting}

\subsection{Part 4: Memory Allocation}

Trace memory allocation:

\begin{lstlisting}[language=bash]
strace -e trace=brk,mmap ./malloc_program
\end{lstlisting}

\subsection{Questions for Reader}

\begin{enumerate}
\item Which system calls does \texttt{printf} use internally?
\item What is the difference between \texttt{write} and \texttt{fwrite}?
\item How does the C library implement buffered I/O?
\item What system calls are used for memory allocation?
\end{enumerate}
